\chapter{Limitations and Conclusion}
\label{ch:conclusion}

\section{Limitations}

\textbf{Model and survey coverage.} All fine-tuning experiments use Qwen3-8B. A
second model family was not evaluated because the Llama-3.1-8B weights are gated
and the account used for this project was not authorized. The pipeline is
model-agnostic, but the empirical claims are conditional on Qwen3-8B. Both ANES
and GSS are from the United States; cross-national surveys, non-English
questionnaires, and non-political attitude domains are not tested.

\textbf{Group attitudes.} A nine-group voter--group corpus and two protocols were
constructed, but the group model was not evaluated. Group coverage is uneven
across waves (three groups in 2016, nine in 2020, two in 2022), and the 2020 block
lacks respondent identifiers. The claims are restricted to candidate-level PAS and
GSS FAS.

\textbf{Panel and choice metric.} ANES is a panel: there are no test-only
respondents in 2022/2024, and the T$\times$R cell is small (226 rows). No protocol
isolates a genuinely new respondent observed only in a future wave. Choice is
rating-induced: it is compared with the highest ground-truth thermometer score,
not with an observed vote.

\textbf{Ablations and FAS.} The ablations use one seed each; the effects are small
and the cross-pressure cells are small. FAS covers two three-attribute
configurations. EDC improves strongly; EHC improves less, and the joint symmetric
KL moves in the opposite direction for direct SFT. The temporal extrapolation is
validated only for 2024, and only one mitigation is tested.

\textbf{Ethics.} Synthetic respondents are not a substitute for real participants.
A model trained on historical data can reproduce and amplify the biases of that
data, and its outputs can be mistaken for evidence about a population. The
temporal results are a warning: a simulator that looks accurate in-distribution
can lose all individual signal four years later. Deployments should report the
training era, validation protocol, model family, and generated-sample coverage,
and synthetic results should be clearly labelled.

\section{Conclusion}

This thesis presented a data-centric pipeline for LLM survey simulation and
evaluated it at the individual, population, and temporal levels. The pipeline
reconstructs interaction-level ANES data, decouples candidate identity from
attributes, encodes partisan and ideological structure as explicit features,
samples the training pool to cover tails and conflict cases, and fine-tunes a
compact 8B model with LoRA. The same pipeline reconstructs a GSS joint
distribution and fine-tunes a model to generate a synthetic population.

The individual-level result is deliberately modest. On unseen respondents in the
same era, the fine-tuned model reaches an MAE of 20.84, a correlation of 0.679,
and a choice accuracy of 0.518 against chance 0.163. A gradient-boosted baseline
reaches 20.58, 0.714, and 0.534, and paired bootstrap shows no significant MAE or
choice difference. The value of the language model is not superior point
prediction; it is a unified interface that produces a rating distribution and
extends to population-level generation. The population-level result is stronger
for EDC: fine-tuning reduces marginal and joint JS divergence by 91\% and 86\%
(93\% and 82\% for the fully covered run).

The temporal result is the most important negative finding. Individual-level
correlation falls from 0.722 in 2019 to $-0.056$ in 2024, and the tabular baseline
collapses in the same way. Adding the most recent wave does not help. Structured
alignment is necessary but not sufficient for temporal generalization, and
temporal shift should be treated as a first-class evaluation axis. Future work
should replicate the pipeline on a second model family, extend it to
cross-national and group-attitude data, develop temporal adaptation methods, and
validate rating-induced choice against observed vote choice.
